\section{Оценка и интерпретация результатов}
Prolog удобен для последовательности вычислений, отрицания как неудачи
и формирования списков кандидатов. OWL удобен для описания общей модели,
наследования, ограничений и проверки непротиворечивости. Эквивалентные классы
позволяют классифицировать индивиды без отдельного запуска каждого правила.
При этом одинаковое внешнее слово «не» имеет различную семантику:
в Prolog достаточно отсутствия доказательства, а в OWL требуется логическое
основание отрицания.

Приложение использует Prolog как источник решений, поскольку условия диалога
естественно выражаются предикатами. Онтология выполняет роль формальной модели
и независимой проверки выбранных классификаций. Совпадение результатов
подтверждено для заданного конечного набора; оно не доказывает эквивалентность
двух формализмов для произвольного расширения базы.

Полнота рекомендаций ограничена пятью записанными жанровыми связями.
Следующий этап развития --- расширение фактов с сохранением их происхождения,
пересмотр явных отрицательных категорий и добавление подтверждённых возрастных
ограничений. Для большого каталога можно использовать постоянный процесс
интерпретатора и отдельное хранилище знаний. До расширения данных такие
изменения не улучшают качество рекомендации сами по себе.

\clearpage\section{Заключение}
Выполнен переход от фактов о видеоиграх к правилам Prolog, онтологии OWL
и диалоговой системе поддержки принятия решений. Количество фактов и правил
соответствует заданию, запросы демонстрируют переменные, логические операции
и вывод. Онтология задаёт иерархию, домены и диапазоны, несовместимость классов,
свойство данных, кардинальности и SWRL-правило.

Тестирование показало совпадение выбранных классификаций Prolog и HermiT,
обнаружение искусственно внесённых противоречий и корректную обработку
положительных и отрицательных сценариев диалога. Пользователь получает
идентификаторы рекомендованных игр и основания их отбора.
Практическая польза подхода состоит в прозрачности: изменение результата
можно проследить до факта, правила или введённого ограничения.

\section{Источники}
\begin{enumerate}[leftmargin=*]
\item Бессмертный И. А., Авдюшина А. Е., Кугаевских А. В., Королева Ю. А.,
Рущенко Н. Г. Системы искусственного интеллекта. Методическое пособие.
СПб.: Университет ИТМО, 2024. Раздел «Базы знаний и онтологии»,
с. 56--58, приложение 2. \url{https://books.ifmo.ru/file/pdf/3308.pdf}.
\item Задания курса «Системы ИИ», файл DOCX, предоставленный к лабораторным.
Модуль 1 и шаблон расширенного отчёта.
\item Образец оформления: отчёт по вычислительной математике,
Numerical-methods/Lab6, репозиторий delovoyhomie/ITMO\_labs.
\end{enumerate}
